Two beliefs dominate how the death overs (16--20) of Twenty20 (T20) cricket are
modelled: that the historical record of the batter and bowler on strike drives the
next ball, and that ``momentum'' or pressure sequences shift outcomes beyond the
game state. We test both at the level of an individual delivery and find little
support for either. On 223{,}678 death-over deliveries from 4{,}402 men's IPL and
T20I matches (Cricsheet, 2005--2026), we build calibrated per-ball models of
expected runs and wicket probability and decompose their skill into game state
versus a historical-rate encoding of identity, using a nested model ladder, grouped
SHAP, and match-clustered bootstrap intervals.

Our headline result concerns momentum. Applying the Miller--Sanjurjo
streak-selection correction to cricket --- to our knowledge for the first time ---
we show the apparent death-over ``cold hand'' (naive $D=-0.114$) is almost
entirely a finite-sequence artifact (bias $-0.121$): the corrected estimate is
$+0.007$ and not significant against a state-stratified null ($p=0.40$). The
$+0.121$ gap between the naive and corrected estimators quantifies the size of the
apparent-momentum illusion. On the identity question, we find no evidence that our
historical-rate encoding of identity improves out-of-sample skill: a
game-state-only model is the best-calibrated wicket model (log-loss $0.2738$ vs.\
$0.2748$ base rate), and adding identity worsens held-out log-loss and
significantly raises expected-runs RMSE in every one of six rolling held-out
seasons --- even though SHAP attributes $51$--$64\%$ of explanatory mass to batter
identity. A random-player placebo collapses that mass to $7$--$8\%$ while skill is
unchanged, showing the attribution is faithful but the signal does not generalise:
a cautionary dissociation between in-sample attribution and out-of-sample skill.
Residual dot-ball pressure adds significant skill to expected runs (RMSE $+0.0065$,
95\% CI $[0.0048, 0.0083]$; longer dot runs lower expected runs) but nothing to
wicket probability, so the Markov property is violated for scoring intensity but
not for dismissal risk. The absolute skill differences are small in proper-score
terms --- these are per-ball outcomes close to the base rate --- so our claims rest
on the \emph{qualitative} dissociation and the bias correction, not on the
magnitude of any single model improvement.
